% ==========================================================================
% Capítulo: Contribución y resultados del proyecto de grado
% Estado: se registra en ../STATUS.md, no aquí.
%
% Debe contener (project-context/formato-anteproyecto.md, "Contribución y
% resultados del proyecto de grado"; extensión máxima 4 páginas):
%   - APORTES RELACIONADOS CON EL OBJETO DEL PROYECTO: impactos respecto a
%     la solución del problema formulado (institucionales, académicos,
%     económicos, de privacidad de datos).
%   - APORTES RELACIONADOS CON EL DESARROLLO DE CAPACIDADES DEL
%     INVESTIGADOR.
%   - RESULTADOS Y ENTREGABLES: por CADA objetivo específico, al menos un
%     resultado con su entregable. La trazabilidad objetivo -> resultado ->
%     entregable debe ser explícita (tabla).
% NO debe contener: beneficios no trazables a
% ../../project-context/documentation.md ("Expected Results",
% "Deliverables", "Benefits").
% ==========================================================================
\chapter{Contribución y resultados del proyecto de grado}\label{cha:contribucion-resultados}

\section{Aportes relacionados con el objeto del proyecto}

\todo{Expresiones como «maximiza el retorno de la inversión» y «reducen la carga cognitiva» (párrafo siguiente) son afirmaciones que el proyecto no medirá. Usen formulaciones como «contribuye a» y retiren o sustenten la referencia a la carga cognitiva.}

La plataforma descrita en los capítulos anteriores incide directamente sobre el problema formulado en el \cref{cha:descripcion-problema}, que reúne la interrupción del ciclo de vida de los modelos tras el entrenamiento y la gestión manual de una infraestructura compartida que no se gobierna ni se observa. Su primer aporte es de naturaleza institucional y económica. Al habilitar el despliegue de modelos ya entrenados sobre el mismo \emph{hardware} que hoy se usa para el entrenamiento, el proyecto maximiza el retorno de la inversión realizada en la infraestructura del laboratorio, en lugar de dejarla limitada a la fase experimental del ciclo de vida. El sistema de reserva de cupos con prioridad por rol, acoplado a los roles institucionales, amplía el acceso a las unidades de procesamiento gráfico entre estudiantes y profesores de la universidad, con reglas de prioridad declaradas para cada rol. Esta misma capacidad reduce la necesidad de adquirir créditos de cómputo en nubes públicas comerciales para sostener cargas de trabajo que la infraestructura \emph{on-premise} del laboratorio puede absorber, argumento que se desarrolla en el \cref{cha:motivacion-antecedentes}.

\todo{La frase sobre «arquitecturas de gran escala, en particular los LLM» que pasan del entrenamiento al servicio de inferencia sugiere que el laboratorio entrena LLM, lo cual no está sustentado, y restringe el aporte a los modelos de lenguaje. Reformúlenla en términos de modelos de IA ya entrenados, propios o descargados.}

En el plano académico, el proyecto habilita la ejecución de proyectos avanzados en áreas prioritarias para el macroproyecto, como Industria 4.0 y e-Salud, al permitir que arquitecturas de gran escala, en particular los LLM, pasen de la fase de entrenamiento a un servicio de inferencia consumible dentro de la misma plataforma institucional. Esta capacidad tiene además una implicación de privacidad y gobernanza de datos, porque al mantener el procesamiento de inferencia sobre servidores propios del IAsLab, en lugar de delegarlo en proveedores externos, la información asociada a estos proyectos de investigación permanece bajo el control administrativo de la universidad. Finalmente, en el plano operativo, el módulo de telemetría y la administración centralizada de cuotas reducen la carga cognitiva y administrativa que hoy recae sobre los administradores del laboratorio, quienes actualmente carecen de visibilidad \emph{fine-grained} sobre el consumo real de la infraestructura para diagnosticar saturaciones o conflictos de uso entre usuarios.

\section{Aportes relacionados con el desarrollo de capacidades del investigador}

\todo{Este párrafo presenta la integración de KubeRay como una competencia central, aunque en el prototipo Ray quedó como una variante comparada. Ajústenlo según lo que decidan en el comentario de la sección 4.4.}

\todo{Revisen que todo elemento que se nombra haya sido presentado antes o tenga una referencia explícita. Por ejemplo, este párrafo menciona «la imagen oficial de \emph{llama.cpp}» como si el lector ya la conociera, y el capítulo 9 usa conceptos como el «fondo común» de cuota que no aparecen en el marco teórico. Hagan esta revisión en todos los capítulos.}

La ejecución de este proyecto exige y desarrolla competencias profesionales específicas de la disciplina, propias de la ingeniería telemática y de sistemas. El despliegue y la configuración del clúster de Kubernetes sobre las estaciones de trabajo del laboratorio, la integración de KubeRay y del NVIDIA GPU Operator, y la definición de reglas de selección de nodos y aislamiento de cargas de trabajo desarrollan competencias de orquestación de contenedores y cómputo distribuido en topologías heterogéneas de CPU y GPU. La integración de motores de inferencia empaquetados como imágenes de contenedor, elegidos por medición y con la imagen oficial de \emph{llama.cpp} como candidato preferente para servir modelos ya cuantizados en formato GGUF dentro de la memoria de video disponible, desarrolla, adicionalmente, competencias de ingeniería de plataformas de inferencia bajo restricciones de \emph{hardware}.

\todo{La sigla CPU (párrafo anterior) se usa sin definir, pues solo se define en el Anexo D, y RBAC se define aquí aunque el control de acceso basado en roles ya aparece en la sección 4.6. Definan cada sigla en su primer uso dentro del texto.}

La implementación del sistema de reserva de cupos con prioridad por rol, acoplado a los atributos de rol expuestos por el servicio de identidad institucional SAAMFI, exige competencias de diseño de esquemas de gobernanza y de integración con sistemas de control de acceso basado en roles (RBAC), un dominio propio de la ingeniería de \emph{software} orientada a sistemas multiusuario. La construcción del módulo de observabilidad, apoyado en Prometheus, Thanos, Grafana, Loki y el exportador DCGM de NVIDIA, desarrolla competencias de ingeniería de observabilidad de infraestructura distribuida, entendida como la capacidad de instrumentar, recolectar y visualizar métricas de \emph{hardware} e inferencia en tiempo real. Por último, la evaluación de la plataforma mediante pruebas de carga sobre los nodos de inferencia, ejecutadas con \emph{scripts} de Python, desarrolla competencias de diseño y ejecución de evaluación empírica de sistemas de \emph{software}, indispensables para sustentar decisiones de ingeniería con evidencia.

\section{Resultados y entregables}

\todo{En la Tabla 8.1, el resultado del objetivo 1 solo incluye métricas de modelos de lenguaje (\emph{tokens} por segundo). Añadan las métricas de los demás modelos, como la latencia y la tasa de peticiones atendidas (R1-24).}

\todo{En la Tabla 8.1, el resultado del objetivo 4 incluye «el crecimiento de la capacidad por nodo», que no forma parte del objetivo tal como está redactado. Alineen la tabla con el objetivo o el objetivo con la tabla.}

De acuerdo con el formato institucional del anteproyecto, cada resultado esperado y su entregable correspondiente deben trazarse a al menos uno de los objetivos específicos formulados en el \cref{cha:objetivos}. La \cref{tab:objetivos-entregables} presenta esta trazabilidad para los cuatro objetivos específicos del proyecto.

\begin{xltabular}{\textwidth}{p{0.18\textwidth} X X}
        \caption{Trazabilidad entre objetivos específicos, resultados esperados y entregables.}
        \label{tab:objetivos-entregables} \\
        \toprule
        \textbf{Objetivo específico} & \textbf{Resultado esperado} & \textbf{Entregable} \\
        \midrule
        \endfirsthead
        \toprule
        \textbf{Objetivo específico} & \textbf{Resultado esperado} & \textbf{Entregable} \\
        \midrule
        \endhead
        \bottomrule
        \endfoot
        \Cref{obj:telemetria}: desarrollo del módulo de observabilidad &
        Panel de telemetría que muestra, por nodo, métricas de \emph{hardware} (temperatura, VRAM, consumo energético) y de carga de trabajo (latencia, \emph{tokens} por segundo, saturación de concurrencia), mediante visualizaciones gráficas interactivas &
        Módulo de observabilidad integrado a la plataforma (Prometheus, Thanos, DCGM Exporter, Grafana, Loki) y su documentación de despliegue \\
        \addlinespace
        \Cref{obj:gobernanza}: implementación del sistema de cuotas por rol &
        Sistema de cuotas y restricciones de infraestructura que reserva cupos de cómputo para las cargas de inferencia, con prioridad por rol, acoplado a los roles provistos por SAAMFI &
        Módulo administrativo de gestión de cuotas integrado a la plataforma web y manual de administración de políticas de cuotas \\
        \addlinespace
        \Cref{obj:orquestacion}: implementación del despliegue de modelos &
        Clúster de cómputo que ejecuta motores de inferencia intercambiables, empaquetados como imágenes de contenedor, sobre el cual un usuario despliega modelos ya entrenados de forma concurrente desde la plataforma &
        Clúster desplegado con su configuración versionada, interfaz de despliegue de modelos integrada a la plataforma y manual de despliegue \\
        \addlinespace
        \Cref{obj:evaluacion}: evaluación de desempeño &
        Reporte de desempeño de la plataforma bajo pruebas de carga con usuarios concurrentes sobre los nodos de inferencia, incluidos los modelos que no son de lenguaje y el crecimiento de la capacidad por nodo &
        Informe de resultados de pruebas de carga, con los datos generados por los \emph{scripts} de prueba de carga \\
\end{xltabular}

Los cuatro entregables específicos de la \cref{tab:objetivos-entregables} convergen en los dos entregables globales exigidos para el proyecto de grado. El primero es el \emph{software} funcional y empaquetado, es decir, el código fuente de la plataforma, desplegado sobre la infraestructura propia del IAsLab y estructurado bajo estándares que permitan su registro como producto tecnológico con derechos de autor. El segundo es la documentación técnica y de usuario, compuesta por los manuales de despliegue de modelos, de visualización de telemetría y de administración de políticas de cuotas que se derivan de cada uno de los módulos anteriores. Ambos entregables globales dependen del cumplimiento conjunto de los cuatro objetivos específicos, porque ningún módulo constituye, por sí solo, la plataforma que exige el objeto del proyecto.
